\scanpage{530}
% E530-01
\mspage{568}
\uline{$\overline{\mathrm{IV}}$}\quad
% E530-02
\deleted{Let}

% E530-03
So let $\mathcal{G}$ be a profinite group, equipped \unsure{with the|with}
% E530-04
\inserted{generators}
$\rho,\sigma\in\mathcal{G}$ satisfying the relations
% E530-05
\[
(1)\qquad \rho^6=\sigma^4=1,\qquad
\rho^3=\sigma^2\quad (\overset{\mathrm{dfn}}{=}\text{``$-1$''}) .
\]
% E530-06
\margin{This simply means that we give ourselves a surjective homomorphism
% E530-07
\diagram{\begin{tikzcd}[ampersand replacement=\&]
\mathrm{SL}(2,\mathbb Z)^{\wedge}\arrow[r] \& \mathcal{G}
\end{tikzcd}}}

% E530-08
A particularly interesting case is the one where
$\rho^3=\sigma^2=1$, to which nevertheless \unsure{I} would not
\unsure{wish} to limit myself. Let $\tau$ be an \unsure{automorphism}
of $\mathcal{G}$ satisfying
% E530-09
\[
(2)\qquad
\tau(\deleted{\unsure{\rho}}\inserted{\sigma})=
\begin{array}{c}\sigma^{-1}\\[-2pt]\inserted{=-\sigma}\end{array},
\qquad \tau(\rho)=\sigma(\rho^{-1})
\]
% E530-10
\margin{Let us introduce \deleted{\unsure{the}} right away the
\unsure{equivalent} element $\tau'=\sigma^{-1}\tau$ satisfying
$\tau'(\rho)=\rho^{-1}$, $\tau'(\sigma)=\sigma^{-1}$
$(\tau'(\varepsilon_0)=\varepsilon_1^{-1},\tau'(\varepsilon_1)=\varepsilon_0^{-1})$.}
% E530-11
(which means that $\tau_\infty$ ``passes to the quotient'' as an
\unsure{automorphism} $\tau$ of $\mathcal{G}$), which implies
% E530-12
\[
(3)\qquad \tau^2=1
\]
% E530-13
\[
(4)\qquad \tau(\varepsilon_0)=\varepsilon_0^{-1},\qquad
\tau(\varepsilon_1)=\varepsilon_1^{-1},\qquad\text{where}
\]
% E530-14
\[
(5)\qquad
\begin{array}{c}
\varepsilon_0=\deleted{\unsure{\alpha}}\sigma\rho\\[-2pt]
\inserted{=-\sigma^{-1}\rho}
\end{array},\qquad
\varepsilon_1=\rho\sigma=\sigma(\varepsilon_0)
=\sigma^{-1}(\varepsilon_0)=\rho(\varepsilon_0).
\]
% E530-15
\margin{\unsure{of the formula}\ \unsure{forgotten}\ \unsure{a} sign $-$,
whereas \unsure{we could} \unsure{from the outset|ask}
simply \unsure{give|take}
$\varepsilon_0=\unsure{\sigma_\infty|\sigma^{-1}}\rho\ldots$}

% E530-16
Put
% E530-17
\[
(5)\qquad \mathcal{G}'=\{1,\tau\}\mathbin{.}\mathcal{G},
\]
% E530-18
so we have
% E530-19
\[
(6)\qquad
\diagram{\begin{tikzcd}[ampersand replacement=\&]
1\arrow[r]\&\mathcal{G}\arrow[r]\&\mathcal{G}'\arrow[r]\&
\{1,\tau\}\arrow[r]\&1.
\end{tikzcd}}
\]

% E530-20
Let
% E530-21
\deleted{\unsure{(7)}}
% E530-22
\[
\begin{gathered}
B_{\rho,\sigma}\subset\operatorname{Aut}(\mathcal{G})\times\widehat{\mathbb Z}^{*}
\\
(7)\qquad =\left\{(u,\mu)\ \left|\
\begin{array}{l}
u(\varepsilon_0)=\varepsilon_0^\mu\\
\deleted{\text{and}}\ u(\sigma)\ \text{conjugate in }\mathcal{G}\text{ to }\sigma^\mu\\
u(\rho)\quad\rule{4cm}{0.3pt}\quad
\mathcal{G}^{\deleted{\unsure{\prime}}}\text{ to }\rho^\mu
\end{array}\right.\right\}.
\end{gathered}
\]
% E530-23
which is obviously a \uline{subgroup}.
The third condition \unsure{can be made explicit} moreover,
% E530-24
\margin{529}

% E530-F01
\margin{\begin{center}
529\\
\texttt{https://grothendieck.umontpellier.fr}\\
Université de Montpellier, UMR CNRS 5149, Case courrier 051, Place Eugène Bataillon\\
34095 Montpellier cedex 5 - France
\end{center}}

\scanpage{531}
% E531-01
since $\rho^6=1$ i.e. $\rho^\mu$ depends only on the image
of $\mu$ in
% E531-02
\[
(\mathbb Z/6\mathbb Z)^*\simeq
\inserted{(\mathbb Z/2\mathbb Z)^*\times}
(\mathbb Z/3\mathbb Z)^*\simeq(\mathbb Z/3\mathbb Z)^*\simeq\pm1,
\]
% E531-03
\deleted{\unsure{only} on \illegible{1}}
% E531-04
\inserted{so}
% E531-05
\[
(8)\qquad \rho^\mu=\rho^{\varepsilon_3(\mu)}
\]
% E531-06
that for $\varepsilon_3(\mu)=1$
% E531-07
\deleted{we}
$u(\rho)$ is conjugate to $\rho$ in $\mathcal{G}$, and for
$\varepsilon_3(\mu)=-1$ it is conjugate to $\rho^{-1}$ in
$\mathcal{G}$, or again to $\rho$ by an element of
$\mathcal{G}'$ -- i.e. an element of the form $\alpha_1\tau_\infty$
with $\alpha_1\in\mathcal{G}$ (since
$\tau_\infty(\rho)=\inserted{\sigma}(\rho^{-1})$ precisely\ldots).
% E531-08
\deleted{We shall thus have,}
% E531-09
\deleted{\[
(9)\qquad u(\rho)=\zeta\rho,\qquad u(\sigma)=\eta\sigma
\quad(\zeta,\eta\in\mathcal{G})
\]}
% E531-10
\deleted{relations}
% E531-11
\deleted{\[
(10)\qquad \zeta=\alpha\rho(\alpha)^{-1},
\qquad\eta=\beta
\]}
% E531-12
We see likewise that we have
% E531-13
\[
\begin{aligned}
(9)\qquad \sigma^\mu=\sigma^{\varepsilon_2(\mu)}
&\quad\left\{\begin{array}{ll}
=\sigma&\text{if }\mu\equiv1\quad(4)\\
\deleted{\unsure{=}}\sigma^{-1}=-\sigma
&\text{if }\mu\equiv-1\quad(4)
\end{array}\right.\\[-2pt]
&\hspace{2.5cm}\deleted{=\unsure{\varepsilon\sigma|\sigma}}\\
&=\varepsilon_2(\mu)\sigma
\end{aligned}
\]
% E531-14
\deleted{But So \unsure{even|=} \unsure{we write}
(for $(u,\mu)\in B_{\rho,\sigma}$) some
\deleted{\unsure{such $\alpha\in$}} elements
$\alpha\in\mathcal{G}'$, $\beta\in$}

% E531-15
Thus:

% E531-16
\uline{Proposition}\quad For \unsure{a|one}
% E531-17
\deleted{have}
$(u,\mu)\in\operatorname{Aut}
(\deleted{\unsure{B_{\rho,\sigma}}}\inserted{\mathcal{G}})
\times\widehat{\mathbb Z}^{*}$ to be in $B_{\rho,\sigma}$,
it is necessary and sufficient that there exist
$\alpha\in\mathcal{G}'_{\deleted{\unsure{\rho,\sigma}}}
{}^{\deleted{\unsure{\varepsilon_3}}}$,
$\beta\in\mathcal{G}$ such that we have
% E531-18
\margin{530}

% E531-F01
\margin{\begin{center}
530\\
\texttt{https://grothendieck.umontpellier.fr}\\
Université de Montpellier, UMR CNRS 5149, Case courrier 051, Place Eugène Bataillon\\
34095 Montpellier cedex 5 - France
\end{center}}

\scanpage{532}
% E532-01
\mspage{569}
% E532-02
\[
(10)\qquad \unsure{u}(\rho)=\operatorname{int}(\alpha)(\rho),\qquad
u(\sigma)=\underbrace{\varepsilon_2(\mu)}_{\varepsilon}
\operatorname{int}(\beta)(\sigma)
\]
% E532-03
and finally
% E532-04
\[
(11)\qquad u(\varepsilon_0)=\varepsilon_0^\mu,
\]
% E532-05
where in addition we have
% E532-06
\deleted{\[
(12)\qquad
\alpha\overset{\inserted{\in\mathcal{G}}}{\unsure{\bmod}\mathcal{G}}
\equiv\tau^{\overset{\inserted{\unsure{\nu_3}}}
{\unsure{\varepsilon_3(\mu)}}}
\quad\overset{\inserted{\mu\equiv1\ (3)}}{(\mathcal{G})}
\quad\text{i.e. }\alpha\in\mathcal{G}\unsure{\Longleftrightarrow}\mu.
\]}
% E532-07
\margin{i.e. $\unsure{\alpha}\equiv\tau$.}

% E532-08
We shall note that
% E532-09
\deleted{$u$}\inserted{$(u,\mu)$}
is uniquely determined by
% E532-10
\inserted{$\mu$ and}
the knowledge of $\alpha$
% E532-11
\inserted{and}
of $\beta$,
% E532-12
\deleted{and of $\varepsilon_2(\mu)$}
and even \unsure{only} by $\mu$ and
% E532-13
\[
(13)\qquad \zeta=\alpha\rho(\alpha)^{-1},\qquad
\eta=\varepsilon_2(\mu)\beta\sigma(\beta)^{-1}
\]
% E532-14
since we \unsure{shall then have} $u$ by
% E532-15
\[
(14)\qquad u(\rho)=\zeta\rho,\qquad u(\sigma)=\eta\sigma.
\]

% E532-16
Moreover, in the case where
% E532-17
\[
(15)\qquad
\diagram{\begin{tikzcd}[ampersand replacement=\&]
\widehat{\mathbb Z}\arrow[r]\&\mathcal{G}\\
n\arrow[r,mapsto]\&\varepsilon_0^n
\end{tikzcd}}
\]
% E532-18
is injective,
% E532-19
\deleted{then}
$\mu$ is determined as a function of $u$ as
% E532-20
\inserted{multiplier of}
the action of
% E532-21
\deleted{\unsure{$u$}}\inserted{$u$}
on $L_0^{\mathcal{G}}\subset\mathcal{G}$, image of (15),
so $(u,\mu)$ is determined by $u$ so by
$(\alpha,\beta,\varepsilon_{\unsure{4}}
\unsure{\Leftarrow|=}\varepsilon_2(\mu))$.
Moreover, if $\mathcal{G}$ is a group that maps into
$\mathrm{SL}(2,\widehat{\mathbb Z})'$
(so as to send $\rho,\sigma$ to the usual
elements), we have seen that $\alpha$ all by itself
% E532-22
\margin{531}

% E532-F01
\margin{\begin{center}
531\\
\texttt{https://grothendieck.umontpellier.fr}\\
Université de Montpellier, UMR CNRS 5149, Case courrier 051, Place Eugène Bataillon\\
34095 Montpellier cedex 5 - France
\end{center}}

\scanpage{533}
% E533-01
already determines $\mu$, so $\varepsilon$, so \deleted{$(\alpha,\beta)$} \deleted{\unsure{determines|all}} the element of $B_{\rho,\sigma}$ is determined with the help of $\alpha,\beta$ alone.

% E533-02
But what conditions must $(\alpha,\beta,\mu)$ satisfy to determine an \unsure{element} of $B_{\rho,\sigma}$? First of all, it \unsure{is necessary} that there indeed exist an \unsure{automorphism} $u$ (necessarily$^{t}$ unique) satisfying (14) — this is a condition on $\zeta,\eta$ alone.

% E533-03
\deleted{It remains to express}
\inserted{We must next express (11). \deleted{\unsure{we}}}
\deleted{Next the relation (11) is written, since
% E533-04
\[
\varepsilon_0=\sigma\rho=\deleted{\unsure{\text{so}}}-\sigma^{-1}\rho,\quad\text{so}
\]
% E533-05
\[
\begin{aligned}
u(\varepsilon_0)&=-u(\sigma^{-1})u(\rho)=-(\eta\sigma)^{-1}\zeta\rho
 =\deleted{\unsure{\sigma(\eta^{-1}\zeta)\underbrace{\sigma\rho}_{\varepsilon_0}}}\\
&\inserted{-\sigma^{-1}\eta^{-1}\zeta\underbrace{\rho\sigma^{-1}}\sigma
 =\sigma^{-1}(\eta^{-1}\zeta\varepsilon_1)}
\end{aligned}
\]
\inserted{in the form}
% E533-06
\[
\eta^{-1}\zeta\varepsilon_1\deleted{\unsure{\sigma(\varepsilon_0)|\sigma\rho}}
 =\varepsilon_0^\mu\quad\text{i.e.}\quad\unsure{\zeta}
\]
% E533-07
\[
\zeta\,\sigma(\eta)=\varepsilon_0^{\mu-1}\;(=l_0^\nu)
\]
% E533-08
\unsure{We shall get from this|We also have}, \deleted{\illegible{1}}
% E533-09
\[
\begin{aligned}
\sigma(\eta)&=\varepsilon\sigma\bigl(\beta\sigma(\beta)^{-1}\bigr)
 =\varepsilon\sigma(\beta)\underbrace{\sigma^2(\beta)^{-1}}_{\inserted{=\beta^{-1}}}
 =\deleted{\unsure{u}}\varepsilon\deleted{\unsure{\sigma}}\bigl(\beta\sigma(\beta)^{-1}\bigr)^{-1}\\
&=\deleted{\unsure{\varepsilon\bigl(\beta\sigma(\beta)^{-1}\bigr)}}\eta^{-1}
\end{aligned}
\]
i.e.
% E533-10
\[
\tag{16}
\deleted{\unsure{\rho}}\quad\eta\sigma(\eta)=\deleted{\unsure{u}}\,1
\]
}

% E533-11
Let us introduce
\[
\tag{*}
\varepsilon'_0=\sigma^{-1}\rho=-\varepsilon_0.
\]
whence
% E533-12
\[
{\varepsilon'_0}^{\mu}=(-1)^\mu\varepsilon_0^\mu=-\varepsilon_0^\mu
\qquad(\text{since }\mu\equiv1\;(2))
\]
% E533-13
\[
u(\varepsilon'_0)=-u(\varepsilon_0)
\]
since $u(-1)=-1$ cf. below.

% E533-14
so the relation (11) is also equivalent to
\[
\tag{\unsure{**|*}}
u(\varepsilon'_0)={\varepsilon'_0}^{\mu},
\]
which, by the same old calculation, is written thus
% E533-15
\margin{532}

% E533-F01
\margin{\begin{center}
532\\
\texttt{https://grothendieck.umontpellier.fr}\\
Université de Montpellier, UMR CNRS 5149, Case courrier 051, Place Eugène Bataillon\\
34095 Montpellier cedex 5 - France
\end{center}}